\section{小元数的补充证明}\label{app:small}
\begin{lemma}\label{lem:triple-large}
三平均在全部 $n\ge11$ 上满足 $G$ 为三幂的充要条件。
\end{lemma}
\begin{proof}
记 $d=d_3(n)$。有 $n\ge d+10$：$d\le1$ 时直接成立，$d=2$ 时 $n\ge14$，$d\ge3$ 时 $n\ge2\cdot5^{d-1}$。引理\ref{lem:safety}及\ref{lem:initialize}的保护池论证适用于三平均，从 $3X$ 至多两步建立两个三重值。
至少四个三重值时选两个同余值平均 $(a,a,b)$，另两重值保留。恰三个时，同余的两种配比至少一个避开第三值；三者各占一个模三类且有轻值 $x$ 时选 $(a,a,x)$，其中 $a\equiv x$。若只有三个实际值，选 $(a,b,c)$：除非仅一值重数超过三，否则保留两重值。余下为 $(a^3,b^3,c^{n-6})$；若均值碰撞到 $c$，则 $a+b=2c$，零和给 $c=0,a=-b$。
仅两个同余三重值时，有富余副本即可用 $(a,a,b)$。否则两者各三份，余池大小至少 $d+4$。若有同余轻值，两个锚点的选择不可能同时碰撞；若余池单一模三类，使用保护池。若余池占另外两类，跨类位置对至少 $n-7>d$ 个，至多 $d$ 个由某个二元危险集阻塞。取未阻塞的 $(c,e)$，$(a,c,e)$ 和 $(b,c,e)$ 安全，不可能同时碰撞到另一锚点。
仅两个不同余三重值时，轻值 $x\equiv a$ 的选择 $(a,a,x)$ 失败只可能 $x=3b-2a$、$f(a)\le4$。若 $x$ 有两份，改用 $(a,x,x)$；未选 $a,b$ 和这一仿射关系保证安全。故两锚点余类至多各有一个额外单点 $x=3b-2a,y=3a-2b$，其余池 $Z$ 在第三余类。
两单点都存在时 $Z\ne\varnothing$，选 $(x,y,z)$。仅 $x$ 存在时也有 $Z\ne\varnothing$，否则安全操作后得到不同余的两值状态，而零和及合法性会迫使 $n\mid f(a)$。选 $(x,b,z)$；$f(b)=3$ 时 $|Z|\ge3$、每值至多两份，可避开 $z=5a-4b$。无额外单点时，两锚点各三份则保护池完成；两者均至少四份则选 $(a,b,z)$。若 $f(a)=3,f(b)\ge4$，选 $z\ne2b-a$。不存在这种 $z$ 时，$Z$ 为一至两份 $2b-a$，零和给 $(3-|Z|)(a-b)+nb=0$，而 $\gcd(a-b,n)=1$，矛盾。
以上安全性均由未选锚点和整数仿射关系给出。非终端步骤严格降低整数平方能量，唯一终端 $(a^3,(-a)^3,0^{n-6})$ 用三次 $(a,-a,0)$ 消去。
\end{proof}

\begin{lemma}\label{lem:eight}
三平均八元的合法配置均可归零。
\end{lemma}
\begin{proof}
若无五重值，存在安全非恒等整数三元组。奇偶分布为 $4+4$ 或 $2+6$；前者任意三元组补集都保留奇偶见证。若没有非恒等整三元组，则至多两个模三类，且每个至少三项的类实际常值；无五重值迫使 $(a^4,b^4)$，由零和得 $b=-a$，违反合法性。在 $2+6$ 情形，若多数池没有整三元组，无五重值迫使它的两模三类各至少两项；其两项和覆盖全部模三类，任一少数位置可与其中两项配成安全整三元组。
整数能量下降至五重值 $u$ 后，写例外为 $a,b,c$，选 $c-u$ 奇。平均 $(u,a,b)$，得到合法 $B_8(u,v)=(u^4,v^3,-4u-3v)$。
令 $x=u,y=u+v$，本原参数满足 $y$ 奇。下列实际返回给三个作用：
\[
R:(x,y)\mapsto(x+y,-y),\quad
S:(x,y)\mapsto(y-x,y),\quad
F:(x,y)\mapsto(3x+3y,2x-4y).
\]
为列出操作，置 $w=-4u-3v$、$a_0=(u+2v)/3$、$b_0=(2u+v)/3$、$q_0=-u-2v/3$、$h_0=-u/3-4v/9$、$k_0=u/3-2v/9$。三个词依次为
\begin{align*}
R:&\ (v,v,u),(v,u,u),(w,a_0,b_0),(a_0,b_0,u);\\
S:&\ (w,v,u),(q_0,q_0,u),(q_0,u,u),
(h_0,v,k_0),(h_0,v,k_0),(v/9,v/9,k_0);\\
F:&\ (v,u,u),(w,b_0,b_0),(v,u,u).
\end{align*}
逐次代入可核对库存及最终重数 $4,3,1$。$SR,RS$ 给平移 $x\mapsto x\pm2y$。若 $|y|>1$，先用 $S$ 使 $x$ 奇，再平移至 $|x_0-2y|<|y|$；等号仅在 $y\mid x$ 时发生，与本原性矛盾。施加 $F$ 后两坐标均偶，约去 $2$ 后第二坐标奇，故本原输出满足 $|y'|\le|x_0-2y|<|y|$。最终 $|y|=1$，平移和 $S$ 送到 $(0,\pm1)$。配置 $(0^4,1^3,-3)$ 先平均 $(-3,0,0)$，再三次平均 $(-1,0,1)$。
\end{proof}

\begin{lemma}\label{lem:ten}
三平均十元的合法配置均可归零。
\end{lemma}
\begin{proof}
先证明无六重值时有安全整三元组。危险奇偶集只能有两项，记为 $P$；危险模五集有两或三项，记为 $F$。若二者均不存在，任意整三元组安全；不存在整三元组会迫使 $(a^5,b^5)$，违反合法性。仅有一种危险集时，在多数池内用模三两类的两项和；若无法配出安全三元组，多数池中至少六项为同一实际值。
两危险集同时存在时，令 $K=(P\cup F)^c$。其内部整三元组安全。若 $|K|\ge7$，无六重值使两个模三类均至少两项，与任一外部位置配对即得安全选择。$|K|=6$ 时只剩 $r^5s$，全部外部余数必须为 $s$；从 $P\setminus F,F\setminus P$ 各取一项，连同 $K$ 中的 $s$ 即可。$|K|=5$ 时 $P,F$ 不交，只剩 $r^4s$ 或 $r^5$；前者同理。后者若 $P,F$ 有跨类余数，配一项 $K$ 得三类各一；否则外部五项同余，取一项 $P$、两项 $F$。每种选择补集均保留两种见证。

能量下降至 $(u^6,a,b,c,d)$。令 $e_a=a-u$ 等，和为 $-10u$。若某个例外差是模 $10$ 单位，平均其余三个例外，一步进入合法 $B_{10}(u,v)=(u^6,v^3,-6u-3v)$。否则四个差模 $10$ 为两份 $5$ 及一对相反非零偶数，故有单位和的跨类配对，设为 $e_a+e_b$。令
\[q=(-5u-a-b)/3,\quad r=(-5u+2a+2b)/9,\quad s=(-11u-a-b)/27.
\]
词 $(u,c,d),(q,a,b),(u,q,r),(u,q,r)$ 给 $u^3r s^6=B_{10}(s,u)$，且 $s-u=-(40u+e_a+e_b)/27$ 是模 $10$ 单位。

在本原 $(X,Y)$ 方向 $X=u,Y=3u+2v$ 上，合法类型为：两者奇且 $X+Y\equiv0\pmod4$，取 $h=1$；或奇偶相反，取 $h=2$。另要求 $5\nmid Y$。整数势能为 $h^2(5X^2+Y^2)$。命题\ref{prop:ten-certificate}给七个实际返回及十二个完整二次覆盖：除 $X=0$ 外，总有返回严格降低该势能。$X=0$ 时已有零坐标，固定一个零，对其余九个位置采用两层三进网络。由整数下降完成。
\end{proof}

\begin{lemma}\label{lem:collision-small}
五平均的 $n=16,17,18$ 及七平均的 $n=23,24,25,26$ 均满足判据。
\end{lemma}
\begin{proof}
引理\ref{lem:initialize}适用。引理\ref{lem:background}、\ref{lem:heavy}、\ref{lem:two-heavy}只留下不同余的两个重值 $a,b$，$f(a)=p+\alpha,f(b)=p+\beta$、$\alpha\ge\beta\ge0$、$\alpha+\beta\le p-3$。五平均十六元另有三值 $(a^5,b^6,c^5)$：一次调节若碰撞，$i=1,3$ 会使三值同奇偶，$i=2$ 给 $b=0,a=-c$，因此也已处理。
从需要保护的轻值危险集选一个包含极小横截集 $T$，$h=|T|\le d_p(n)$。每个 $v\in T$ 有私人素数 $\ell$，其危险支持与 $T$ 恰交于 $v$；此处 $a=b\pmod\ell$，其余保护位置与 $b$ 同余。故未保护部分相对 $b$ 的总和不可能为 $0\pmod\ell$。
令 $\Delta=a-b$、$t_i=(x_i-b)/\Delta$、$z_i=t_i\pmod p$。锚点同余修复后，余数 $0,1$ 各至多一个，实际参数为 $p,1-p$；各余类重数小于 $p$。对轻值子集 $S$，$k=|S|\le p$，若 $i\equiv-\sum_{S}z\pmod p$ 满足 $i\le\min(\alpha,p-k)$，碰撞要求
\begin{equation}\label{eq:collision}
\sum_{S}t=p-i;
\end{equation}
若还满足 $p-k-i\le\beta$，两个旧重值均保留，因而该模式不可能失败。从另一锚点，$j\equiv\sum_S z-k\pmod p$、$j\le\min(\beta,p-k)$ 时，碰撞要求 $\sum_S t=k+j-p$。
命题\ref{prop:collision-certificate}穷尽这些线性系统。五平均令 $c=n-10-h\in\{5,6,7,8\}$。$c=7,8$ 无解；$c=6$ 的两个解均为整数且 $D=5+\alpha+\sum t_i=15$；$c=5$ 时所有 $D$ 的分母为奇数。对 $n=17,18$，有保护时私人素数整除 $\Delta$，而 $D=15$ 给总和矛盾；无保护时没有相容模式。对 $n=16,h=0$，$16b+15\Delta=0$ 迫使 $\Delta$ 偶，全部整数参数使配置同奇偶；$h=1$ 时 $\Delta$ 偶而 $D$ 分母奇，同样违反私人素数 $2$ 的总和条件。
七平均可在每个余数类内对\eqref{eq:collision}取对称平均，压缩为变量 $y_r$、重数 $f_r$；总和保持。某类在一个方程中部分选取时，交换位置进一步强迫该类实际常值。令 $n-h=21+e$，则 $1\le e\le5$。证书给十一种相容模式，每个实际参数均为整数，且 $0<D=7+\alpha+\sum f_ry_r<21+e$。$h>0$ 时私人素数给总和矛盾；$h=0$ 时所有值属于 $b+\Delta\Z$，合法性给 $\gcd(\Delta,n)=1$，零和迫使 $n\mid D$，又矛盾。因此总有安全续步，整数能量下降完成。
\end{proof}

\begin{lemma}\label{lem:seven22}
七平均二十二元的合法配置均可归零。
\end{lemma}
\begin{proof}
入口使用命题\ref{prop:entry}证明中转换前的 $(a^{14},(-2a+z)^7,-7z)$。三条三步返回为
\[
A=\begin{pmatrix}4&1\\-1&0\end{pmatrix},\quad
F=\begin{pmatrix}1&2\\0&7\end{pmatrix},\quad
J=\begin{pmatrix}1&2\\2&-1\end{pmatrix}.
\]
$A$ 保留一份旧 $A$、平均两组 $(A^6,B)$ 再平均 $(A,B^5,C)$；$F$ 保留 $C$、平均两组 $(A^5,B^2)$ 再平均 $(A^4,B^3)$；$J$ 保留一份旧 $B$、平均两组 $(A^5,B^2)$ 再平均 $(A^4,B^2,C)$。实际尺度均为 $1/7$。有 $(FA^2)^2=-7I,J^2=5I$，故各逆可实现。
令 $H=AFJA^{-1}=\left(\begin{smallmatrix}-7&-62\\0&5\end{smallmatrix}\right)$，则 $[F,H]=U(396/35)$。倍率生成 $R=\Z[1/35]$；用 $C=\left(\begin{smallmatrix}1&4\\0&-1\end{smallmatrix}\right)$ 换基，$CU(t)C=U(-t)$、$CAU(t)A^{-1}C=L(t)$，得到 $\Gamma(396^2,R)$。
记 $M=1936$、$X=F^2/7$。命题\ref{prop:seven-certificate}给 $A^{440}=X^{440}=I\pmod M$，且两者模 $9$ 生成全部 $\SL_2(\Z/9\Z)$。其中三个方向 $I+3E_{12},I+3E_{21},I+3\diag(1,-1)$ 的立方逐层生成模 $81$ 的高层核。因此对任意 $g\in\Gamma(M,R)$，可选可实现 $w\equiv I\pmod M$、$w\equiv g\pmod{81}$；由 $396^2=81M$ 得 $w^{-1}g$ 已在深核中，故 $\Gamma(M,R)$ 可实现。
同一证书给 $R^\times$ 模 $M$ 的像 $U$ 有 $440$ 个元素，全部单位为 $U\sqcup3U$。合法列 $(z,y)=(z,a+7z)$ 的有限标记是 $(z/y,\epsilon)$，$\epsilon$ 记录 $y$ 所属单位陪集。证书证明全部 $3872$ 个标记都能到 $(0,0)$。引理\ref{lem:fiber}再精确送到 $(0,s)$，其中 $s\in R^\times$。配置 $(s^{14},(-2s)^7,0)$ 中，四份 $s$、两份 $-2s$ 和一个零组成零和七元组，接引理\ref{lem:zero-tail}。
\end{proof}

